\documentclass[10pt,a4paper]{amsart}
\usepackage[margin=19mm]{geometry}
\usepackage{amsmath,amssymb,mathtools}
\usepackage[scr=rsfs,cal=euler]{mathalfa}
\usepackage{xfrac}
\usepackage[unicode,hidelinks,bookmarks=false]{hyperref}
\newcommand{\ie}{i.e.\ }
\newcommand{\cf}{cf.\ }
\newcommand{\resp}{respectively\ }
\newcommand{\Subsection}{\subsection}
\providecommand{\nobreakdash}{\nobreak\hbox{-}}
\setlength{\emergencystretch}{2em}
\multlinegap0pt
\usepackage{enumerate}
\usepackage{amssymb}
\usepackage[shortlabels]{enumitem}
\usepackage{dsfont}
\usepackage{tikz}
\newtheorem{theorem}{Theorem}
\title{Two problems of SI-compact spaces}
\author{Zhengmao He}
\date{Source: arXiv:2608.02020v1 (CC BY 4.0)}
\begin{document}
\maketitle

\noindent\emph{Source and licence.} Two problems of SI-compact spaces. Version arXiv:2608.02020v1 (CC BY 4.0), distributed by arXiv under the Creative Commons Attribution 4.0 International licence, http://creativecommons.org/licenses/by/4.0/, as recorded in the arXiv licence metadata for this exact version and shown on the abstract page https://arxiv.org/abs/2608.02020v1. Full licence notice: \texttt{licenses/arxiv-2608.02020-CC-BY-4.0.txt}.\par\smallskip

\noindent\textit{Editorial scope.} Counted unit: the article's theorem theorem (source0.tex lines 262--263). The article's complete original proof of it is delivered with it (source0.tex lines 264--277, one \texttt{proof} environment of 14 lines). No mathematics is written, repaired or re-derived: the statement and the proof are the article's own lines, in the article's own notation, and no retained line is substituted or re-flowed.  No further source-local context is needed: every cross-reference of every retained line resolves inside the two delivered spans.  Nothing was omitted from any retained span, so the package carries no omission record.  No retained line contains a citation, so the wrapper carries no bibliography and no bibliography entry.  No retained line contains a figure environment, an image-inclusion command or drawing code, so no image is delivered and the package declares no asset dependency.  Editorial formatting only, in addition: an AMS \texttt{amsart} preamble that restates the article's own declarations and macros used by the retained lines, namely the environments \texttt{theorem} on separate counters, together with the standard packages those lines need.  The retained environments print in this document as Theorem 1 in order of appearance, whereas the article prints the same items with the numbering of its own full document.  The complete native source of the article as distributed in the arXiv e-print for this exact version is bundled unchanged as \texttt{sources/206659/source0.tex}.
\begin{theorem}{\rm Let $X,Y$ be $SI$-compact spaces. Then $X\times Y$ is $SI$-compact.}
\end{theorem}

\begin{proof} Let $\mathcal{U}\subseteq\mathcal{O}(X\times Y)$ be a Scott irreducible open cover of $X\times Y$. Define a mapping $$\xi:\Sigma\mathcal{O}(X\times Y)\longrightarrow\Sigma\mathcal{O}(X)$$ as follows:
$$\forall\ W\in\mathcal{O}(X\times Y),\ \xi(W)=\{x\in X\mid \{x\}\times Y\subseteq W\}.$$

$\mathbf{Claim} \ 1$: $\xi$ is continuous.

Let $E\in\mathcal{O}(X\times Y)$ and $x\in \xi(E)$. Equivalently, $\{x\}\times Y\subseteq E$. Since $Y$ is $SI$-compact, $Y$ is compact. By Lemma 3.4, there is an open set $F\in\mathcal{O}(X)$ such that $\{x\}\times Y\subseteq F\times Y\subseteq E$. In other words, $x\in F\subseteq\xi(E)$. So $\xi(E)$ is open in $X$ and thus $\xi$ is well defined. Suppose $\{G_{i}\mid i\in I\}\subseteq\mathcal{O}(X\times Y)$ is a directed family. Obviously, $\bigcup\limits_{i\in I}\xi(G_{i})\subseteq \xi(\bigcup\limits_{i\in I}G_{i})$. Let $a\in\xi(\bigcup\limits_{i\in I}G_{i})$. Then by Lemma 3,3, $\{e_{a}(G_{i})\mid i\in I\}$ is an open cover of $Y$. By the compactness of $Y$, there are $G_{i_{1}},G_{i_{2}},\cdot\cdot\cdot,G_{i_{n}}$ such that $Y=e_{a}(G_{i_{1}})\cup e_{a}(G_{i_{2}})\cup\cdot\cdot\cdot\cup e_{a}(G_{i_{n}})$. This means that $\{a\}\times Y\subseteq G_{i_{1}}\cup G_{i_{2}}\cup\cdot\cdot\cdot\cup G_{i_{n}}$. Using the directedness of the family $\{G_{i}\mid i\in I\}$, we can choose a $G_{i^{\ast}}$ such that $G_{i_{1}},G_{i_{2}},\cdot\cdot\cdot,G_{i_{n}}\subseteq G_{i^{\ast}}$. Consequently, $\{a\}\times Y\subseteq G_{i^{\ast}}$ and hence $a\in\xi(G_{i^{\ast}})$. Thus, $\bigcup\limits_{i\in I}\xi(G_{i})=\xi(\bigcup\limits_{i\in I}G_{i})$. Therefore, $\xi$ is continuous.


$\mathbf{Claim}\ 2$: for every $x\in X$, there is a $U_{x}\in\mathcal{U}$ such that $x\in \xi(U_{x})$

By Lemma 3.3, $e_{x}(\mathcal{U})$ is irreducible in $\Sigma\mathcal{O}(X)$. Furthermore, $$\bigcup e_{x}(\mathcal{U})=\bigcup\limits_{U\in\mathcal{U}}e_{x}(U)=Y$$ because $\bigcup\mathcal{U}=X\times Y$. By the $SI$-compactness of $Y$, there is a $U_{x}\in\mathcal{U}$ satisfying $e_{x}(U_{x})=Y$. Whence, $\{x\}\times Y\subseteq U_{x}$. Thus, $x\in\xi(U_{x})$.

By Claim 2, $X\subseteq\bigcup\limits_{x\in X}\xi(U_{x})\subseteq\bigcup\xi(\mathcal{U})=\bigcup\limits_{U\in\mathcal{U}}\xi(U)$. So $\xi(\mathcal{U})$ is an open cover of $X$. By Claim 1, $\xi(\mathcal{U})$ is irreducible in $\Sigma\mathcal{O}(X)$. As $X$ is $SI$-compact, $X\in\xi(\mathcal{U})$. So there is a $U^{\ast}\in\mathcal{U}$ such that $X=\xi(U^{\ast})$. As a consequence, $X\times Y=U^{\ast}\in\mathcal{U}$. Therefore, $X\times Y$ is $SI$-compact.
\end{proof}

\end{document}
